\documentclass[11pt,a4paper]{article}
\usepackage[margin=19mm,headheight=15pt]{geometry}
\usepackage{fontspec}
\setmainfont{Noto Serif}
\usepackage{amsmath,amssymb,mathtools}
\usepackage{xcolor}
\usepackage{fancyhdr}
\usepackage{enumitem}
\usepackage{microtype}
\setlength{\parindent}{0pt}
\setlength{\parskip}{6pt}
\definecolor{navy}{HTML}{183153}
\definecolor{lightgray}{HTML}{F3F5F7}
\pagestyle{fancy}\fancyhf{}
\lhead{\small ME 5309 -- Mid-Semester Examination}
\rhead{\small ME 5309}
\cfoot{\thepage}
\begin{document}
\begin{center}
{\Large\bfseries ME 5309 -- Mid-Semester Examination}\\[3pt]
{\large Answer Booklet}\\[8pt]
\end{center}
\textbf{Instructions:} Attempt any three questions. Each question carries 10 marks. Show working where appropriate.
\section*{Question 2 \hfill [10 marks]}
\subsection*{Q2(a) \hfill [2 marks]}

At equilibrium the velocity should be zero.
The torque and gravity torque must balance.

\subsection*{Q2(b) \hfill [2 marks]}

\[\begin{aligned}A = [[0,1], [-61.3125,-2.5]], B = [[0],[50]]\end{aligned}\]

\subsection*{Q2(c) \hfill [4 marks]}

The balance is 50(90-θss) = 2.4525 sin θss.
Solving gives a small error of about 0.05°.

\subsection*{Q2(d) \hfill [2 marks]}

The added gravity torque should cancel the load, so the steady-state error becomes zero. 
\section*{Question 3 \hfill [10 marks]}
\subsection*{Q3(a) \hfill [2 marks]}

Finite differencing amplifies encoder noise, making the velocity estimate noisy.

\subsection*{Q3(b) \hfill [2 marks]}

I use ωc=20 rad/s and α=ωcT/2=0.01.
Then:\[\begin{aligned}y[k] = 0.980198y[k-1] + 0.009901(x_k+x[k-1])\end{aligned}\]
The DC gain is approximately 1.

\subsection*{Q3(c) \hfill [4 marks]}

\[\begin{aligned}L = [18,64]^T\end{aligned}\]

\subsection*{Q3(d) \hfill [2 marks]}

An observer can estimate velocity from position measurements.
A filter can be safer because it smooths noise.

\section*{Question 4 \hfill [10 marks]}
\subsection*{Q4(a) \hfill [2 marks]}

The closed-loop poles are -2±2j.
I use Nff=c/b=6 for unity DC gain.

\subsection*{Q4(b) \hfill [2 marks]}

The nominal controller has unity DC gain, so θss=1 and the tracking error is zero even after the actuator gain changes.

\subsection*{Q4(c) \hfill [4 marks]}

I obtain s³+(a+bk₂)s²+(c+bk₁)s-bK\_I.
Desired polynomial is s³+8s²+24s+32, so k₂=6, k₁=18, K\_I=-32.

\subsection*{Q4(d) \hfill [2 marks]}

Integral action removes steady-state error, although this only holds for the nominal values of b and c.
Saturation can cause windup, which should be handled by anti-windup.

\end{document}